\begin{thebibliography}{99}
\bibitem{artendinglamination} J. F. Brock, R. D. Canary \& Y. N. Minsky - “The classification of Kleinian surface groups, II: The ending lamination conjecture” , Ann. of Math. (2) 176 (2012) no. 1, p. 1-149
\bibitem{livremartine} M. Babillot, R. Feres \& A. Zeghib - Rigidité, groupe fondamental et dynamique , Panoramas \& Synthèses , vol. 13 , Société Mathématique de France, Paris, 2002
\bibitem{artbishopjonescriticalgeometricfiniteness} C. J. Bishop \& P. W. Jones - “Hausdorff dimension and Kleinian groups” , Acta Math. 179 (1997) no. 1, p. 1-39
\bibitem{articlebishopjones} C. J. Bishop \& P. W. Jones - “The law of the iterated logarithm for Kleinian groups” , in Lipa’s legacy (New York, 1995) , Contemp. Math. , vol. 211 , American Mathematical Society, Providence, RI, 1997, p. 17-50
\bibitem{thesemoi} A. Boulanger - Quelques exemples de systèmes dynamiques: comptage en mesure infinie, enlacement sur le tore et échanges d’intervalles affines , Ph. D. Thesis, Sorbonne Université, 2018
\bibitem{artmoiheatcounting1} A. Boulanger - “A stochastic approach to counting problems” , Ann. Sci. École Norm. Sup. (4) 55 (2022) no. 1, p. 225-259
\bibitem{articlecanary} R. D. Canary - “Ends of hyperbolic $3$-manifolds” , J. Amer. Math. Soc. 6 (1993) no. 1, p. 1-35
\bibitem{surveycanarymarden} R. D. Canary - “Marden’s tameness conjecture: history and applications” , in Geometry, analysis and topology of discrete groups , Adv. Lect. Math. (ALM) , vol. 6 , Int. Press, Somerville, MA, 2008, p. 137-162
\bibitem{arttamenesscalegarigabai} D. Calegari \& D. Gabai - “Shrinkwrapping and the taming of hyperbolic 3-manifolds” , J. Amer. Math. Soc. 19 (2006) no. 2, p. 385-446
\bibitem{livrechavel} I. Chavel - Eigenvalues in Riemannian geometry , Pure and applied math. , vol. 115 , Academic Press, 1984
\bibitem{artdelsarte} J. Delsarte - “Sur le gitter fuchsien” , C. R. Acad. Sci. Paris 214 (1942), p. 147-179
\bibitem{artdodziuk} J. Dodziuk - “Maximum principle for parabolic inequalities and the heat flow on open manifolds” , Indiana Univ. Math. J. 32 (1983) no. 5, p. 703-716
\bibitem{arteskinmac} A. Eskin \& C. McMullen - “Mixing, counting, and equidistribution in Lie groups” , Duke Math. J. 71 (1993) no. 1, p. 181-209
\bibitem{articleepsteinspectraltheory1} C. L. Epstein - The spectral theory of geometrically periodic hyperbolic $3$-manifolds , Mem. Amer. Math. Soc. , vol. 58, no. 335 , American Mathematical Society, Providence, RI, 1985
\bibitem{articleepsteinspectraltheory2} C. L. Epstein - “Asymptotics for closed geodesics in a homology class, the finite volume case” , Duke Math. J. 55 (1987) no. 4, p. 717-757
\bibitem{articleepsteinspectraltheory3} C. L. Epstein - “Positive harmonic functions on abelian covers” , J. Funct. Anal. 82 (1989) no. 2, p. 303-315
\bibitem{artgrigoheatkernelhyper} A. Grigor’yan \& M. Noguchi - “The heat kernel on hyperbolic space” , Bull. London Math. Soc. 30 (1998) no. 6, p. 643-650
\bibitem{bookgrigorheatkernel} A. Grigor’yan - Heat kernel and analysis on manifolds , AMS/IP Studies in Advanced Math. , vol. 47 , American Mathematical Society, Providence, RI, 2009
\bibitem{arthuber2} H. Huber - “Über eine neue Klasse automorpher Funktionen und ein Gitterpunktproblem in der hyperbolischen Ebene. I” , Comment. Math. Helv. 30 (1956), p. 20-62
\bibitem{articlelax} P. D. Lax \& R. S. Phillips - “The asymptotic distribution of lattice points in Euclidean and non-Euclidean spaces” , J. Funct. Anal. 46 (1982) no. 3, p. 280-350
\bibitem{articlesullivandemicylindre} - Geometry Symposium, Utrecht 1980 , Lect. Notes in Math. 894 (1981)
\bibitem{artmargulismixing} G. A. Margulis - “Certain applications of ergodic theory to the investigation of manifolds of negative curvature” , Funkcional. Anal. i Priložen. 3 (1969) no. 4, p. 89-90
\bibitem{livremcmumu} C. T. McMullen - Renormalization and 3-manifolds which fiber over the circle , Annals of Math. Studies , vol. 142 , Princeton University Press, Princeton, NJ, 1996
\bibitem{surveyminsky} Y. N. Minsky - “End Invariants and the classification of hyperbolic 3-manifolds” , Current Developments in Math. (2002), p. 111-141
\bibitem{livreotal2001hyperbolization} J.-P. Otal - The hyperbolization theorem for fibered 3-manifolds , SMF/AMS Texts and Monographs , vol. 7 , American Mathematical Society, Providence, RI; Société Mathématique de France, Paris, 2001, Translated from the 1996 French original (Astérisque, vol. 235)
\bibitem{artpattersontermeerror} S. J. Patterson - “A lattice-point problem in hyperbolic space” , Mathematika 22 (1975) no. 1, p. 81-88
\bibitem{arttersondensiteconforme} S. J. Patterson - “The limit set of a Fuchsian group” , Acta Math. 136 (1976) no. 3-4, p. 241-273
\bibitem{articlepatteron2} S. J. Patterson - “On a lattice-point problem in hyperbolic space and related questions in spectral theory” , Ark. Mat. 26 (1988) no. 1, p. 167-172
\bibitem{artroblinfoncorb} T. Roblin - “Sur la fonction orbitale des groupes discrets en courbure négative” , Ann. Inst. Fourier (Grenoble) 52 (2002) no. 1, p. 145-151
\bibitem{memoireroblin} T. Roblin - Ergodicité et équidistribution en courbure négative , Mém. Soc. Math. France (N.S.) , vol. 95 , Société Mathématique de France, Paris, 2003
\bibitem{artselberg2} A. Selberg - “Harmonic analysis and discontinuous groups in weakly symmetric Riemannian spaces with applications to Dirichlet series” , J. Indian Math. Soc. (N.S.) 20 (1956), p. 47-87
\bibitem{artselberg} A. Selberg - “On discontinuous groups in higher-dimensional symmetric spaces” , in Contributions to function theory (internat. Colloq. Function Theory, Bombay, 1960) , Tata Institute of Fundamental Research, Bombay, 1960, p. 147-164
\bibitem{artsullivandensiteconforme} D. Sullivan - “The density at infinity of a discrete group of hyperbolic motions” , Publ. Math. Inst. Hautes Études Sci. 50 (1979), p. 171-202
\bibitem{artsullivanmesurebms} D. Sullivan - “Entropy, Hausdorff measures old and new, and limit sets of geometrically finite Kleinian groups” , Acta Math. 153 (1984) no. 3-4, p. 259-277
\bibitem{articlesullivanpositivity} D. Sullivan - “Related aspects of positivity in Riemannian geometry” , J. Differential Geom. 25 (1987) no. 3, p. 327-351
\end{thebibliography}
